\chapter{Pterygium}
\index{Pterygium}
\label{sec:Pterygium}

\section{Overview}
Pterygium is a wing-shaped fibrovascular overgrowth of the bulbar conjunctiva that extends onto the cornea. It is associated with chronic ultraviolet (UV) light exposure and is more common in individuals living in sunny or tropical climates. The growth can cause irritation, astigmatism, and in advanced cases, visual obstruction.

\section{Classification}
\begin{description}[font=\normalfont\bfseries,leftmargin=3em,labelwidth=2.5em]
  \item[Cause] Chronic UV light exposure with subsequent conjunctival elastotic degeneration
  \item[Organ System] Eyes
  \item[Pathophysiology] UV radiation damages limbal stem cells and conjunctival epithelium, triggering aberrant fibroblast proliferation, angiogenesis, and extracellular matrix deposition. Chronic inflammation and upregulation of growth factors (VEGF, TGF-$\beta$) drive progressive fibrovascular growth.
  \item[Duration] Chronic and progressive
  \item[Onset] Gradual
  \item[Mode of Transmission] Not transmissible
  \item[Clinical Course] Slow progression over years; may remain stationary or advance onto the cornea toward the visual axis.
  \item[Severity] Mild, moderate, or advanced depending on corneal extent
  \item[Extent of Disease] Localized (unilateral or bilateral)
  \item[Age Group] More common in adults aged 20--50 years
  \item[Epidemiology] Prevalence ranges from 3--20\% in tropical regions; more common in males. Strong association with outdoor occupations (farmers, fishermen, construction workers).
  \item[ICD-11 Code] 9A03.0
\end{description}

\section{Causes}
Chronic exposure to ultraviolet B radiation from sunlight is the primary causative factor. Additional contributory factors include chronic ocular surface irritation from dust, wind, and dryness. Genetic predisposition and human papillomavirus infection have been proposed but remain unconfirmed.

\section{Risk Factors}
\begin{itemize}[noitemsep]
  \item Prolonged UV light exposure (outdoor occupation or lifestyle)
  \item Living near the equator or at high altitude
  \item Male sex
  \item Advanced age
  \item Chronic ocular surface irritation (dust, wind, sand)
  \item Family history of pterygium
\end{itemize}

\section{Signs and Symptoms}
\begin{itemize}[noitemsep]
  \item Fleshy, triangular growth on the conjunctiva extending onto the cornea
  \item Ocular irritation, redness, and foreign body sensation
  \item Blurred vision or distorted vision from induced astigmatism
  \item Tearing and dryness
  \item Difficulty performing daily activities
\end{itemize}

\section{Progression and Stages}
Pterygium typically begins as a small area of conjunctival thickening at the nasal limbus (primary pterygium). It may remain static or slowly advance across the cornea. When the leading edge reaches or approaches the visual axis, significant astigmatism and visual impairment occur. Post-surgical recurrence is common (20--50\%) if adjunctive therapy is not used.

\section{Complications}
\begin{itemize}[noitemsep]
  \item Corneal astigmatism and visual impairment
  \item Restriction of extraocular movement (large or recurrent lesions)
  \item Post-surgical recurrence with more aggressive growth
  \item Chronic ocular surface inflammation
  \item Reduced quality of life
  \item Emotional and psychological effects
\end{itemize}

\section{Diagnosis}
\begin{itemize}[noitemsep]
  \item Slit-lamp examination confirming fibrovascular conjunctival growth onto cornea
  \item Measurement of corneal extension and visual axis involvement
  \item Keratometry or corneal topography for astigmatism assessment
  \item Lab tests (not needed)
  \item Imaging tests (anterior segment OCT for detailed assessment)
\end{itemize}

\section{Prevention}
\begin{itemize}[noitemsep]
  \item Wear UV-blocking sunglasses and a wide-brimmed hat outdoors
  \item Use lubricating eye drops in dry or dusty environments
  \item Avoid prolonged unprotected sun exposure
  \item Regular check-ups
\end{itemize}

\section{Treatment Options}
\begin{itemize}[noitemsep]
  \item Lubricating eye drops and mild anti-inflammatory drops for symptom relief
  \item Sunglasses and environmental protection
  \item Surgical excision with conjunctival autograft or amniotic membrane graft
  \item Adjunctive mitomycin C or intraoperative beta irradiation to reduce recurrence
  \item Lifestyle changes
\end{itemize}

\section{Living with It}
\begin{itemize}[noitemsep]
  \item Follow treatment plan
  \item Regular appointments
  \item Adjust daily habits
  \item Support groups
  \item Keep learning
\end{itemize}

\section{Outlook}
Small pterygia can be managed conservatively without affecting vision. Surgical excision is effective for symptomatic or visually significant lesions, though recurrence remains a challenge. With modern surgical techniques (conjunctival autograft), recurrence rates are below 10--15\%. Most patients maintain good visual function after appropriate management.
